\section{Isolating and prescribing eigenvalue crossings}\label{n:sec-angular}

The angular construction prepares a finite family of modes whose average
frequencies are smaller than \(k\).  It does this by moving eigenlines
through different positions in the ordered spectrum.  There are two steps.
First, we construct a metric with an isolated double eigenvalue at each
prescribed pair of consecutive positions in a long spectral band.  We then
join short paths through these doubles into one closed path.  A continued
eigenline exchanges its ordered position at a double, so the return
permutation of the closed path can be made into a cycle on the whole band.
The frequency estimate will follow by averaging the ordered positions
visited by each line.
The first lemmas isolate doubles and preserve them while other
multiplicities split. Proposition~\ref{n:prop-angular-program} then
packages the closed path, its continued eigenlines, and the strict
mean-frequency margin used in the radial construction.

All metrics in this section are smooth metrics on \(S^2\).  Write \(g_0\)
for the unit round metric and
\[
 d\mu=\frac{d\operatorname{vol}_{g_0}}{4\pi}.
\]
We use the nonpositive convention for the Laplacian.  The eigenvalues of
\(-\Delta_h\), repeated according to multiplicity, are numbered from \(1\);
in particular \(\lambda_1(h)=0\).  For the round metric, put
\begin{equation}\label{n:eq-round-blocks}
 B_l=(l+1)^2,\qquad B_{-1}=0,\qquad
 \lambda_{B_{l-1}+1}(g_0)=\cdots=\lambda_{B_l}(g_0)=l(l+1).
\end{equation}
The corresponding real eigenspace is denoted by \(\mathcal Y_l\).  It
consists of the restrictions of degree-\(l\) homogeneous harmonic
polynomials and has dimension \(2l+1\).

Fix a number \(\vartheta>1\), independently of the large integer \(k\),
and set
\begin{equation}\label{n:eq-band}
 \begin{gathered}
 M=\lfloor\vartheta k\rfloor,\qquad L=\lfloor\sqrt{k}\rfloor,
 \qquad p=B_M=(M+1)^2,\\
 I=\{B_L+1,\ldots,p\},\qquad N=|I|=p-B_L.
 \end{gathered}
\end{equation}
We will cycle the whole round blocks of degrees \(L+1,\ldots,M\),
leaving the lower positions fixed.  The lower cutoff tends to infinity,
which permits crossings between consecutive degree blocks, but is
\(o(k)\), so leaving these positions fixed does not affect the leading band
average.  The lemmas below are stated for general finite cutoffs before
being applied to this band.

\paragraph{The Euclidean comparison space.}
The preceding dimensions give \eqref{n:euclidean-count}.
For completeness, the polynomial characterization follows from the
mean-value property. A smooth radial averaging kernel at scale $R$
reproduces a Euclidean harmonic function. Moving any $k+1$ derivatives
onto the kernel bounds that derivative at a fixed point by $O(R^{-1})$
under the degree-$k$ growth assumption. Letting $R$ tend to infinity
makes all such derivatives zero. Each homogeneous part of the resulting
polynomial is harmonic. In three variables the dimension in degree $l$
is $\binom{l+2}{2}-\binom{l}{2}=2l+1$, with the second term zero for
$l<2$: the polynomial Laplacian onto degree $l-2$ is surjective because
its adjoint in the monomial inner product
$\langle x^\alpha,x^\alpha\rangle=\alpha!$ is multiplication by
$|x|^2$, which is injective.

\subsection{A fixed angular measure and the splitting formula}

For a smooth real function \(w\), define
\begin{equation}\label{n:eq-conformal-family}
 A(w)=\int_{S^2}e^{2w}\,d\mu,\qquad
 \rho_w=\frac{e^{2w}}{A(w)},\qquad
 h^0(w)=\rho_w g_0.
\end{equation}
The metric \(h^0(w)\) has area \(4\pi\), and its normalized area form is
\(d\mu_w=\rho_w\,d\mu\).  We work in a small convex neighborhood
\(\mathcal U\) of \(0\) among smooth functions, with smallness measured in
\(C^2\).  Given any prescribed \(0<\kappa_*<1\), it can be chosen
so that \(K_{h^0(w)}>\kappa_*\) and so that the
pointwise comparison of \(h^0(w)\) with \(g_0\) is as close to equality as
needed below.  Every path and perturbation used in the construction will
belong to a finite-dimensional smooth family contained in \(\mathcal U\).

The radial argument will use one Hilbert space \(L^2(d\mu)\), so we also
fix coordinates in which the angular area form is independent of \(w\).
Here is a concrete version of Moser's volume-form construction
\cite[Section~4, Theorem~2]{Moser}.  For \(0\leq\tau\leq1\), set
\[
 \rho_{\tau,w}=1+\tau(\rho_w-1).
\]
This is positive and has integral \(1\).  Solve, with mean zero,
\[
 \Delta_0\phi_{\tau,w}=-\partial_\tau\rho_{\tau,w},
 \qquad
 X_{\tau,w}=\rho_{\tau,w}^{-1}\nabla_0\phi_{\tau,w}.
\]
The right side of the Poisson equation has mean zero, and hence the
solution exists uniquely.  If \(\Phi_{\tau,w}\) is the flow of
\(X_{\tau,w}\), starting at the identity, then
\[
 \frac{d}{d\tau}\Phi_{\tau,w}^*(\rho_{\tau,w}\,d\mu)
 =\Phi_{\tau,w}^*
   \left[\bigl(\partial_\tau\rho_{\tau,w}
          +\operatorname{div}_0(\rho_{\tau,w}X_{\tau,w})\bigr)d\mu\right]=0.
\]
Put \(\Phi_w=\Phi_{1,w}\) and
\begin{equation}\label{n:eq-fixed-area-gauge}
 h(w)=\Phi_w^*h^0(w).
 \qquad\text{Then}\qquad
 d\operatorname{vol}_{h(w)}=d\operatorname{vol}_{g_0}=4\pi\,d\mu .
\end{equation}
The Poisson solution and flow depend smoothly on every additional smooth
finite parameter.  The construction depends only on \(w\), and
\(\Phi_0=\operatorname{id}\).  Thus a closed path of \(w\)'s gives a
closed path of \(h(w)\)'s.  On each compact finite parameter family the
metrics in this gauge have bounded smooth geometry.  Their spectra are
those of \(h^0(w)\), and their curvature remains greater than \(\kappa_*\).
If \(w\) is antipodally even, the Poisson solution is even, its vector
field is equivariant under the antipodal map, and its flow commutes with
that map.  Both forms
of the metric then preserve the even and odd function spaces.

The gauge also preserves uniform closeness to the round metric.  To see
this without a bound on the higher derivatives of \(w\), fix
\(0<\alpha<1\) and put \(\sigma=\|w\|_{C^2}\).  For small \(\sigma\),
normalization gives
\[
 \|\rho_w-1\|_{C^1}+\|\rho_w-1\|_{C^{0,\alpha}}\leq C\sigma.
\]
The mean-zero Poisson solution above is independent of \(\tau\), since
its right side is \(-(\rho_w-1)\).  The elliptic estimate on the fixed
round sphere bounds its \(C^{2,\alpha}\) norm by \(C\sigma\), and hence
\(\|X_{\tau,w}\|_{C^1}\leq C\sigma\), uniformly for \(0\leq\tau\leq1\).
The variational equation for the flow gives
\[
 e^{-C\sigma}g_0\leq\Phi_w^*g_0\leq e^{C\sigma}g_0.
\]
Multiplication by \(\rho_w\circ\Phi_w\) gives the same bounds for
\(h(w)\), after increasing \(C\).  Consequently, for every
\(0<\delta<1\), a sufficiently small convex \(C^2\) neighborhood
\(\mathcal U\) ensures
\begin{equation}\label{n:eq-gauge-closeness}
 (1-\delta)g_0\leq h(w)\leq(1+\delta)g_0
 \qquad(w\in\mathcal U).
\end{equation}
This choice is uniform over all finite parameter families contained in
\(\mathcal U\).

The following variation formula is the local input for constructing and
avoiding coincidences.  Its round-sphere full-block version is part of the
classical conformal spectral transversality theory of Colin de Verdi\`ere
and Greilhuber--Kepplinger
\cite{ColinDeVerdiere,GreilhuberKepplinger}.  We give the two-dimensional
and parity-restricted statements that will be used here.

\begin{lemma}[Two independent splitting directions]\label{n:lem-two-directions}
Let \(\lambda>0\) be an eigenvalue of \(h^0(w)\), and let \(u,v\) be an
orthonormal pair in its real eigenspace, using \(L^2(d\mu_w)\).
The traceless compression to \(\operatorname{span}\{u,v\}\) of the
first variation of the Laplacian has rank two as the smooth conformal
direction varies.  If \(w\) is antipodally even and \(u,v\)
have the same antipodal parity, the conclusion still holds when only
even conformal directions are allowed.  The same splitting matrices apply
after the fixed-area coordinate change \eqref{n:eq-fixed-area-gauge}.
\end{lemma}

\begin{proof}
In dimension two the normalized Dirichlet form of \(h^0(w)\) is
\[
 q(f,g)=\int_{S^2}\langle\nabla_0 f,\nabla_0g\rangle\,d\mu,
\]
independent of \(w\).  Its mass form is
\(m_w(f,g)=\int f g\,d\mu_w\).  Differentiating the latter gives
\[
 \dot m_w(f,g)=
 2\int_{S^2}fg
    \left(\dot w-\int_{S^2}\dot w\,d\mu_w\right)d\mu_w.
\]
For an \(m_w\)-orthonormal basis \(u_1,\ldots,u_r\) of a repeated
positive eigenspace, differentiation of the generalized eigenvalue
equation \(q=\lambda m_w\) therefore gives the symmetric cluster matrix
\begin{equation}\label{n:eq-conformal-variation}
 \dot M_{ij}=
 -2\lambda\int_{S^2}u_i u_j
       \left(\dot w-\int_{S^2}\dot w\,d\mu_w\right)d\mu_w .
\end{equation}
Terms arising from a change of orthonormal frame cancel at a scalar
cluster matrix.  The mean term in \eqref{n:eq-conformal-variation} is
scalar, so the two traceless entries on the pair \(u,v\) are represented,
up to fixed nonzero factors, by
\begin{equation}\label{n:eq-product-pair}
 U_1=u^2-v^2,\qquad U_2=2uv.
\end{equation}
Both have mean zero for \(d\mu_w\).

The functions \(U_1,U_2\) are linearly independent.  Indeed, a nonzero
linear combination of them is a nonzero traceless quadratic form in
\((u,v)\); its zero set in \(\mathbb R^2\) is the union of two lines.
If that quadratic form vanished at every point of the sphere, choose a
point where \((u,v)\ne(0,0)\) and a connected neighborhood on which this
remains true.  Continuity then puts \((u,v)\) in just one of the two
lines on a smaller neighborhood.  A fixed nonzero linear combination of
\(u,v\) would vanish there.  Unique continuation for smooth Laplace
eigenfunctions on the connected sphere forces it to vanish everywhere
\cite[\S5, Remark~3]{AKS}, contradicting orthonormality.  Independence
of \(U_1,U_2\) makes their Gram matrix for \(d\mu_w\) positive definite.
Testing in the directions \(U_1,U_2\), for example, proves rank two.
When \(u,v\) have the same parity, both functions in
\eqref{n:eq-product-pair} are even, so the same tests are available in
the even family.

Finally, conjugating the Laplacian by the parameter-dependent pullback
adds to its variation a commutator with the Laplacian.  The compression
of such a commutator to an eigenspace on which the Laplacian is
\(\lambda I\) is zero.  Thus the coordinate change does not change the
splitting matrix.  The zero eigenvalue requires no splitting argument:
it is simple because \(S^2\) is connected.
\end{proof}

\subsection{Avoiding unwanted multiplicities}

We say that a metric is \emph{simple through position \(q\)} when
\[
 \lambda_1<\lambda_2<\cdots<\lambda_q<\lambda_{q+1}.
\]
The final inequality is part of the definition.  At a metric that fails
this condition, the repeated cluster meeting the cutoff may continue
past \(q+1\); the argument below always includes that entire cluster.
For an antipodally invariant metric we use the same convention for the
ordered even and odd lists separately, with any prescribed finite cutoff
in each list.

\begin{lemma}[Finite-cutoff avoidance]\label{n:lem-avoidance}
Parameters whose metrics are simple through any fixed position are dense
in \(\mathcal U\).
Any two such parameters can be joined inside \(\mathcal U\) by a smooth
path that is simple through that position and is constant near its
endpoints.  In the even parameter family, the analogous statements hold
for simplicity through prescribed finite cutoffs within each parity.
Cross-parity coincidences are allowed in this last statement.
\end{lemma}

\begin{proof}
We give the finite-dimensional reduction, including the treatment of
clusters of multiplicity greater than two.  The use of spectral
projections and local cluster matrices is the usual perturbation
construction \cite[III, \S6.4, Theorem~6.17; IV, \S3.4,
Theorem~3.16]{Kato}. In the fixed-area gauge the operators have common
domain $H^2(S^2)$ in $L^2(\mu)$ and depend smoothly, as maps
$H^2\to L^2$, on the finite parameters. On a contour in the resolvent
set, elliptic estimates give bounded inverses $L^2\to H^2$;
the inverse-difference identity and a Neumann series give smooth
parameter dependence. Integrating the resolvents around the contour
therefore gives a smooth finite-rank projection. This argument needs
smooth parameter dependence, rather than analytic eigenvalue branches.
The avoidance step is a finite-dimensional
general-position argument of the type used in spectral genericity
\cite{Uhlenbeck}.

Start with a smooth preliminary path \(w_0(t)\), \(0\leq t\leq1\),
joining the two parameters and constant near them.  Convexity of
\(\mathcal U\) supplies such a path.  The endpoints have positive gaps
through the cutoff, so a small initial and final time interval remains
simple under all sufficiently small perturbations of the path.  Choose
a smooth cutoff \(\chi(t)\) that vanishes near \(0,1\) and equals \(1\)
outside these two simple intervals.

Fix the cutoff \(q\).  Uniform comparison with the round mass form
bounds \(\lambda_{q+1}(h^0(w))\) above throughout a small neighborhood
of the path.  Choose a larger number \(\Lambda\).  Every repeated
cluster that can affect simplicity through \(q\) has eigenvalue below
\(\Lambda\), including all members of a cluster that straddles the
cutoff.  The reverse form comparison bounds the number of eigenvalues
below \(\Lambda\) uniformly.  Consider all triples
\[
 (t,u,v),\qquad
 -\Delta_{h^0(w_0(t))}u=\lambda u,\quad
 -\Delta_{h^0(w_0(t))}v=\lambda v,\quad 0<\lambda\leq\Lambda,
\]
where \(u,v\) are real and orthonormal for \(d\mu_{w_0(t)}\).
This set of normalized pairs is compact in the smooth topology.
To see the needed compactness directly, write the eigenvalue equation
as \(-\Delta_0u=\lambda\rho_{w_0(t)}u\).  The coefficients have uniform
smooth bounds along the compact path, the mass norms are uniformly
equivalent, and \(\lambda\) is bounded.  Elliptic estimates followed by
Sobolev embedding give bounds for every spatial derivative.  A diagonal
subsequence then converges smoothly, and the equation and
orthonormality pass to the limit.  The first positive eigenvalue is
uniformly bounded away from zero, so the limit is again a positive
eigenvalue pair.  In the even construction take only pairs of the same
specified parity; this is a closed condition.
If there is no such pair, the preliminary path is already simple through
the cutoff and no perturbation is needed.

At each such pair, Lemma~\ref{n:lem-two-directions} supplies two smooth
directions for which the two traceless entries have an invertible
derivative.  This property is open in \((t,u,v)\).  A finite cover of
the compact set of pairs therefore supplies a single finite list of
directions \(\xi_1,\ldots,\xi_D\) whose span \(E\) has rank two on
every one of these pairs.  In the parity version all \(\xi_j\) are
even.  The same rank assertion holds for relevant eigenfunction pairs
at all parameters in a sufficiently small \(E\)-neighborhood of the
path.  Otherwise there would be a sequence of counterexamples with
parameters tending to the path.  The same elliptic compactness would
give a limiting pair on the path, where one of the selected two-by-two
determinants is nonzero, a contradiction.

Use this fixed pool in the finite family
\begin{equation}\label{n:eq-path-perturbation}
 w(t,z)=w_0(t)+\chi(t)\sum_{j=1}^D z_j\xi_j,
 \qquad z=(z_1,\ldots,z_D).
\end{equation}
For small \(z\) it stays in \(\mathcal U\); the endpoint intervals
remain simple.  At every possible bad point outside those intervals,
\(\chi=1\), and the two traceless entries have rank two in the
\(z\) variables.

We now describe precisely the sets to be avoided.  At a bad point
\((t_*,z_*)\), let \(r,\ldots,s\) be the \emph{entire} consecutive
block occupied by one repeated eigenvalue, with strict spectral gaps
immediately outside that block.  This includes \(s>q+1\) if the
cutoff meets a larger cluster.  A contour enclosing this eigenvalue
and no other spectrum gives, in a neighborhood of \((t_*,z_*)\), a
smooth Riesz projection of rank \(m=s-r+1\).  Projecting a basis at
the center and orthonormalizing gives a smooth real frame; elliptic
regularity makes it smooth in the spatial variables as well.  In this
frame the operator on the cluster is a smooth symmetric matrix
\(M(t,z)\).  Choose the first two frame vectors at the center to be
any orthonormal pair in the repeated eigenspace and set
\begin{equation}\label{n:eq-chart-equations}
 F(t,z)=\bigl(M_{11}(t,z)-M_{22}(t,z),\,2M_{12}(t,z)\bigr).
\end{equation}
At the center the derivative of \(F\) in \(z\) has rank two, by the
choice of \(E\) and \eqref{n:eq-conformal-variation}.  Shrinking the
neighborhood keeps this rank.  The implicit-function theorem makes
\(F^{-1}(0)\) there a smooth codimension-two submanifold of time
times parameter space.

If at a point in this neighborhood the same block \(r,\ldots,s\) is
still entirely equal, then \(M\) is scalar and \(F=0\).  This is the
only containment asserted for this chart.  When an \(m\)-fold block
partly splits but leaves a smaller repeated block, the smaller block
has its own chart centered at that point.  Thus the two equations
in \eqref{n:eq-chart-equations} handle every \(m\geq2\): they contain
the locus of an exact \(m\)-fold block without purporting to describe
all partial repetitions of the larger Riesz cluster.

For each fixed pair of indices \(r,s\), take the stratum on which
\(r,\ldots,s\) is the entire repeated block, with strict gaps outside
it.  This stratum, as a subspace of the second-countable space
\([0,1]\times\mathbb R^D\), has a countable subcover by neighborhoods
centered on that same stratum.  Every point of the stratum in one of
these neighborhoods lies in that chart's zero set, because the same
block is entirely equal there.  Taking the union over the countably
many possible pairs \(r,s\) covers all bad points by countably many
codimension-two zero sets.  Each zero set projects to a null set in
\(\mathbb R^D\).  For example, in a local implicit-function chart
two coordinates of \(z\) are smooth functions of \(t\) and the
remaining \(D-2\) coordinates.  Its projection is locally a smooth,
hence locally Lipschitz, image of a space of dimension \(D-1\) in
\(\mathbb R^D\), and has \(D\)-dimensional Lebesgue measure zero.
A countable union still has measure zero.  There are therefore
arbitrarily small choices of \(z\) whose path meets none of the bad
sets.  Formula \eqref{n:eq-path-perturbation} fixes the endpoints.

The endpoint density assertion is the same argument without the time
variable.  Begin at one parameter, choose the finite pool for its
normalized low-eigenvalue pairs, and keep its rank in a small
neighborhood.  Each exact-block zero set then has codimension two in
the finite parameter space itself, so the union of the bad sets is
null and its complement is dense.  In the even family, perform all
of these constructions separately for the two parity lists with the
common finite pool of even directions.  The union of their bad sets
is still countable and null.  This proves both versions.
\end{proof}

\subsection{Keeping one double while separating the rest}

Avoidance supplies simple paths, but we also need deliberately chosen
doubles.  The next lemma allows one isolated double to remain while
other coincidences are removed.  Its proof uses only a nonscalar
variation of each unwanted cluster within the two linearized
constraints for the chosen double.  We do not need a description of
the full multiplicity locus of that unwanted cluster.

\begin{lemma}[Retaining an isolated double]\label{n:lem-retained-double}
Suppose that \(h^0(w_*)\) has an isolated eigenvalue of exact
multiplicity two at positions \(i,i+1\).  For any \(q\geq i+1\) and
any neighborhood of \(w_*\) in \(\mathcal U\), there is a parameter
in that neighborhood for which the same positions are equal and all
other inequalities through position \(q\), including the gap after
\(q\), are strict.  The perturbations here are unrestricted conformal
perturbations, even when \(w_*\) was obtained in the even family.
\end{lemma}

\begin{proof}
Let \(u,v\) be an orthonormal basis for the chosen double at the
current parameter, and put \(U_1=u^2-v^2\), \(U_2=2uv\).
The linearized constraints for keeping its cluster matrix scalar
have kernel
\begin{equation}\label{n:eq-double-kernel}
 \mathcal K=\left\{\xi\in C^\infty(S^2;\mathbb R):
       \int U_1\xi\,d\mu_w=\int U_2\xi\,d\mu_w=0\right\}.
\end{equation}
This has codimension two in the full space of smooth conformal
directions, by Lemma~\ref{n:lem-two-directions}.  We first prove that,
for any different repeated positive eigenspace, some direction in
\(\mathcal K\) has a nonscalar compression to that eigenspace.

Suppose instead that every \(\xi\in\mathcal K\) had scalar
compression on such an eigenspace, of eigenvalue \(\mu_*>0\).
Choose an orthonormal pair \(a,b\) in it.  By
\eqref{n:eq-conformal-variation}, the two independent functionals
represented by
\[
 A_1=a^2-b^2,\qquad A_2=2ab
\]
would both vanish on \(\mathcal K\).  A linear functional that
vanishes on the kernel of the surjection
\(\xi\mapsto(\int U_1\xi\,d\mu_w,\int U_2\xi\,d\mu_w)\)
is a linear combination of those two component functionals.
Since the two \(A\)-functionals are themselves independent, their
span must equal the span of the two \(U\)-functionals.  Equality of
the functionals on every smooth \(\xi\), and positivity of
\(\rho_w\), then gives the pointwise identity
\begin{equation}\label{n:eq-quadratic-relation}
 Q(a,b)=TQ(u,v),\qquad
 Q(x,y)=(x^2-y^2,\,2xy),
\end{equation}
for a constant invertible real \(2\times2\) matrix \(T\).
There is no undetermined additive constant here: each of the four
product functions has mean zero.  This argument took place in the
full smooth tangent space before any finite-dimensional restriction.

Identity \eqref{n:eq-quadratic-relation} forces an equality of nodal
sets.  In fact \(a=0\) is equivalent to
\(Q(a,b)\in\{(-t,0):t\geq0\}\).  The inverse image of this closed
ray under \(T\) is another closed ray through the origin.  Identifying
\(\mathbb R^2\) with \(\mathbb C\), the map \(Q\) is the squaring map,
so the inverse image under \(Q\) of a closed ray is one unoriented
line.  There is consequently a nonzero pair \((c,d)\) such that
\begin{equation}\label{n:eq-common-nodal-set}
 \{a=0\}=\{cu+dv=0\}.
\end{equation}
The assertion includes points where both entries of either pair
vanish, because the rays include the origin.

Here is why \eqref{n:eq-common-nodal-set} is impossible for distinct
eigenvalues.  Write \(f=cu+dv\), a nonzero eigenfunction for the
chosen eigenvalue \(\lambda_*\), and take a connected component
\(\Omega\) of the common nonzero set.  Change the signs so that both
\(a\) and \(f\) are positive on \(\Omega\).  Each restriction belongs
to \(H^1_0(\Omega)\), without any regularity assumption on the nodal
boundary.  For example,
\[
 f_\varepsilon=(f-\varepsilon)_+
\]
has compact support in \(\Omega\): its support stays away from the
boundary, where \(f=0\).  As \(\varepsilon\downarrow0\), these
truncations converge to \(f\) in \(H^1(\Omega)\), by dominated
convergence for the functions and their gradients.  Each truncation
can be approximated in \(H^1\) by smooth functions compactly
supported in \(\Omega\).  The same reasoning applies to \(a\).
Their weak eigenvalue equations may therefore be tested against one
another on \(\Omega\).  Symmetry of the Dirichlet form gives
\[
 (\mu_*-\lambda_*)\int_\Omega af\,d\mu_w=0.
\]
The integral is positive.  The two eigenvalues would be equal,
contrary to the choice of a different spectral cluster.  This proves
the existence of a nonscalar direction \(\xi\in\mathcal K\).

We now realize that tangent direction while keeping the chosen double
exactly.  Pick two smooth directions \(\zeta_1,\zeta_2\) on which
the derivative of its two traceless entries is invertible.  In the
three-parameter family
\[
 w(t,x_1,x_2)=w+t\xi+x_1\zeta_1+x_2\zeta_2
\]
form the smooth \(2\times2\) Riesz cluster matrix \(M_{\rm d}\) of
the chosen double and its two constraints
\(F_{\rm d}=(M_{{\rm d},11}-M_{{\rm d},22},2M_{{\rm d},12})\).
They vanish exactly when that cluster is scalar.  The derivative
in \((x_1,x_2)\) is invertible, so the finite-dimensional
implicit-function theorem gives smooth \(x_1(t),x_2(t)\) with
\(F_{\rm d}=0\).  Because \(\xi\in\mathcal K\), their derivatives
at \(0\) are zero.  The resulting curve has tangent \(\xi\).
On the other repeated cluster its matrix is therefore
\[
 \mu_* I+tD+o(t)
\]
with \(D\) nonscalar.  For a sufficiently small nonzero \(t\),
the distinct eigenvalues of \(D\) separate that cluster into at
least two smaller groups.  The chosen cluster remains a double.

Only finitely many repetitions must be removed.  Initially extend
the list past \(q+1\) to the last member \(q'\) of the cluster
containing \(\lambda_{q+1}\).  There is a strict gap after \(q'\).
At each step take the curve parameter small enough to preserve that
gap, the isolation of the chosen double, and every strict gap already
obtained among these \(q'\) positions.  If the multiplicities of
the current clusters are \(m_1,\ldots,m_r\), the integer
\(\sum_j\binom{m_j}{2}\) strictly decreases whenever an unwanted
cluster is split, while the contribution of the chosen double stays
equal to \(1\).  Thus finitely many steps separate every other
cluster.  The steps may be chosen with arbitrarily small total size,
and the isolation gaps keep the chosen double at positions \(i,i+1\).
The claimed neighborhood and cutoff conclusions follow.
\end{proof}

\subsection{Pole-regular modes of a rotational sphere}

The cross-block construction will use the lowest mode in a fixed azimuthal
sector and a specified position in the axisymmetric spectrum. We record
their simplicity and parity directly, including the pole conditions.

\begin{lemma}[Sector simplicity and reflection parity]
\label{n:lem-pole-modes}
Let $\zeta(\theta)>0$ be smooth as an axisymmetric function on $S^2$,
and suppose $\zeta(\pi-\theta)=\zeta(\theta)$. In the metric
$\zeta^2g_0$, fix a nonnegative integer azimuthal number $m$ and one
azimuthal factor (the constant factor if $m=0$). Every eigenvalue of the
meridional profile problem is simple. Its lowest profile can be chosen
strictly positive on $(0,\pi)$ and is even under equatorial reflection.
Along a smooth reflection-symmetric deformation, every fixed position
in the $m=0$ list retains its reflection parity. In particular the
position occupied at round by degree $l$ retains parity $(-1)^l$.
\end{lemma}

\begin{proof}
The separated equation for a profile $v$ is
\[
 -(\sin\theta\,v')'+\frac{m^2}{\sin\theta}v
   =\lambda\zeta^2\sin\theta\,v.
\]
Its form domain is the closure of smooth sphere sector profiles in the
energy-plus-mass norm. Compact spectral theory on the smooth compact
sphere, restricted to this closed sector, supplies its eigenfunctions
and the variational minimum. Their profiles are smooth at the poles:
for $m\geq1$ they have the form $t^m a(t^2)$ in polar distance $t$,
and for $m=0$ they are smooth even functions of $t$.
For two profiles $u,v$ at the same eigenvalue, the equation makes
$W=\sin\theta(uv'-u'v)$ constant. The stated pole behavior makes
$W\to0$ at a pole: for $m\geq1$, $u,v=O(t^m)$ and their derivatives
are $O(t^{m-1})$; for $m=0$, the functions are bounded and their
derivatives are $O(t)$. Thus $W=0$. At any interior point where $u$
is nonzero, $v$ and a scalar multiple of $u$ have equal value and
derivative. Interior ODE uniqueness makes them equal everywhere.

Replacing a minimizing real profile by its absolute value preserves the
form domain and its energy and mass. The resulting nonnegative minimizer
satisfies the weak equation and is smooth in the open interval. An
interior zero would be a minimum with zero derivative, so ODE uniqueness
would force the profile to vanish. The ground profile is consequently
positive. Reflection preserves the equation and pole conditions. Simplicity
and positivity force the mass-normalized ground profile to be reflection
even. With its azimuthal factor its antipodal parity is $(-1)^m$.

For $m=0$, simplicity holds at every position by the same Wronskian
argument. Min--max gives continuous ordered eigenvalues along a smooth
deformation, and a local simple spectral projection gives a continuous
unit eigenfunction. Reflection acts on its line by $+1$ or $-1$; this
sign is continuous and hence constant. At round it is the sign of the
corresponding Legendre polynomial, $(-1)^l$. No nodal count for a
singular interval problem is required.
\end{proof}

\subsection{An isolated double at every adjacent band position}

Recall the band \(I\) in \eqref{n:eq-band}, consisting of the whole
round blocks of degrees \(L+1,\ldots,M\).  An adjacent pair in it
either lies inside one round block or straddles the boundary of two
consecutive blocks.
The first case uses full control of a round block's splitting matrix.
The second uses even metrics to force one parity value past a value
of the other parity.

\begin{lemma}[Adjacent doubles]\label{n:lem-adjacent-doubles}
Fix \(\vartheta>1\).  For every prescribed sufficiently small
neighborhood \(\mathcal U\) as above, every sufficiently large integer
\(k\), and every \(i\) with \(B_L+1\leq i<p\), there is
\(w\in\mathcal U\) for which
the only failure of simplicity through position \(p+1\) is
\(\lambda_i(h(w))=\lambda_{i+1}(h(w))\).  This eigenvalue has exact
multiplicity two.  In particular
\(\lambda_p(h(w))<\lambda_{p+1}(h(w))\).
\end{lemma}

\begin{proof}
\emph{Pairs inside one round block.}
Suppose that \(i,i+1\) belong to the round degree-\(l\) block.
Multiplication of harmonic functions defines
\begin{equation}\label{n:eq-symmetric-square}
 \operatorname{Sym}^2\mathcal Y_l
 \longrightarrow\bigoplus_{j=0}^l\mathcal Y_{2j},
 \qquad u\mathbin{\odot}v\longmapsto uv.
\end{equation}
The target contains the products because a homogeneous polynomial of
degree \(2l\) decomposes, on the sphere, into harmonic components of
degrees \(2l,2l-2,\ldots,0\).  The map is an isomorphism.  This is the
round-sphere product-space fact in
Greilhuber--Kepplinger \cite[Proposition~4.6]{GreilhuberKepplinger};
the following argument records the particular algebra that we need.

Its image is invariant under rotations, hence under the rotation
Casimir \(\Delta_0\).  It contains
\[
 F_l(z)=(1-z^2)^l
       =\bigl(\operatorname{Re}((x+\mathrm i y)^l)\bigr)^2
        +\bigl(\operatorname{Im}((x+\mathrm i y)^l)\bigr)^2
       \quad\text{on }S^2 .
\]
For functions of \(z\),
\(\Delta_0=(1-z^2)\partial_z^2-2z\partial_z\), and direct
differentiation gives
\begin{equation}\label{n:eq-casimir-recursion}
  \bigl(\Delta_0+2l(2l+1)\bigr)F_l=4l^2F_{l-1}.
\end{equation}
Expand \(F_l=\sum_{j=0}^l c_{lj}P_{2j}(z)\) in Legendre
polynomials.  Its top coefficient is nonzero by polynomial degree,
and for \(j<l\), \eqref{n:eq-casimir-recursion} gives
\[
 c_{lj}=
 \frac{4l^2}{2l(2l+1)-2j(2j+1)}\,c_{l-1,j}.
\]
Starting with \(F_0=1\), every \(c_{lj}\) is therefore nonzero.
Polynomial spectral projections in \(\Delta_0\) put a nonzero zonal
vector of each \(\mathcal Y_{2j}\) in the image of
\eqref{n:eq-symmetric-square}.  Rotations of a nonzero zonal harmonic
span its harmonic space: they are the harmonic projections of powers
of linear forms, and such powers span the homogeneous polynomials.
The map is consequently onto.  Its source and target both have
dimension
\[
 \frac{(2l+1)(2l+2)}2
  =\sum_{j=0}^l(4j+1)=(l+1)(2l+1),
\]
so it is an isomorphism.

Dualizing this isomorphism in the variation formula
\eqref{n:eq-conformal-variation} shows that the trace-free
first-variation map for the entire round degree-\(l\) cluster is onto
the trace-free symmetric matrices on \(\mathcal Y_l\).  More
explicitly, a nonzero trace-free matrix \(C\) gives the nonzero
product function \(\sum C_{\alpha\beta}u_\alpha u_\beta\) by
injectivity of \eqref{n:eq-symmetric-square}, so it cannot annihilate
every conformal direction.  Choose finitely many directions on which
this derivative is onto, and represent the cluster by its smooth
Riesz matrix in that finite family.  The submersion theorem realizes
every sufficiently small trace-free matrix as its trace-free part.
Choose a diagonal one whose ordered entries are strictly increasing
except for the desired adjacent equality.  It may be made
arbitrarily small.  The round gaps to the other degree blocks remain
positive, so this is an isolated double at the required positions.
Lemma~\ref{n:lem-retained-double}, with cutoff \(p+1\), separates all
other repetitions while retaining it.

\emph{Pairs across consecutive round blocks.}
These pairs have \(i=B_l\) with \(L+1\leq l\leq M-1\).
Put \(P=(-1)^l\) and \(c_0=3/2\).  Choose once and for all a smooth
axisymmetric function \(\psi(\theta)\), supported in an equatorial
band whose closure is disjoint from the poles, where \(\theta\) is
colatitude, such that
\begin{equation}\label{n:eq-equatorial-bump}
 \psi(\pi-\theta)=\psi(\theta),\qquad
 0\leq\psi\leq1,\qquad \psi(\pi/2)=1,\qquad
 \bar\psi=\frac1\pi\int_0^\pi\psi(\theta)\,d\theta<\frac1{10}.
\end{equation}
For large \(l\), consider the unnormalized metric
\begin{equation}\label{n:eq-parity-metric}
 \widehat h_l=\zeta_l(\theta)^2g_0,\qquad
 \zeta_l=1-\frac{c_0}{l}\psi.
\end{equation}
It is antipodally invariant and tends to \(g_0\) in every fixed
smooth norm.  It corresponds to the conformal parameter
\(w_l=\log\zeta_l\).  Dividing the metric by its area factor
multiplies all its eigenvalues by the same positive number, so it
does not change any ordering used below.

The Dirichlet form is unchanged and
\((1-c_0/l)^2\leq\zeta_l^2\leq1\).  Min--max on either parity
subspace therefore places the positions belonging at round to
degree \(q\) in the interval
\begin{equation}\label{n:eq-parity-form-bounds}
 \bigl[q(q+1),\,(1-c_0/l)^{-2}q(q+1)\bigr].
\end{equation}
The intervals for \(q\) and \(q+2\) in the same parity are disjoint
through \(q=l+1\), for all large \(l\).  Indeed the gap between their
round values is \(4q+6\), while, uniformly for \(q\leq l+1\),
\[
 \bigl((1-c_0/l)^{-2}-1\bigr)q(q+1)
 \leq (2c_0+o(1))(q+1)<4q+6.
\]
Here \(c_0<2\) gives the last inequality uniformly for large \(l\).
Thus the indicated degree groups keep their order inside each
parity list.

We next find a value from the parity-\(P\) degree-\(l\) group above
a value from the parity-\(-P\) degree-\((l+1)\) group.  The sector
min--max statements use the Friedrichs form domains obtained by
closing smooth sphere sector profiles in their energy-plus-mass norms.
Near either pole, a smooth sphere sector profile of azimuthal number \(m\) has the form
\(t^{|m|}a(t^2)\), where \(t\) is distance to that pole and \(a\) is
smooth; for \(m=0\), these are smooth profiles even at the poles.
For the first value, take the lowest eigenvalue in azimuthal sector \(m=l\) of
\(\widehat h_l\).  By Lemma~\ref{n:lem-pole-modes}, the meridional ground profile is
positive and unique up to scale.  Reflection in the equator preserves it; the
azimuthal factor acquires \((-1)^l\) under the antipodal map.  Hence
the resulting mode has parity \(P\).  Choose a real eigenfunction
\(u_l\) with \(\int u_l^2\,d\mu=1\), and write its round energy and
changed mass as
\[
 E_l=\int_{S^2}|\nabla_0u_l|^2\,d\mu,\qquad
 M_l=\int_{S^2}\zeta_l^2u_l^2\,d\mu.
\]
The azimuthal term gives
\[
 E_l\geq l^2\int_{S^2}\frac{u_l^2}{\sin^2\theta}\,d\mu,
 \qquad M_l\leq1.
\]
Min--max in this sector bounds
\[
 \frac{E_l}{M_l}\leq(1-c_0/l)^{-2}l(l+1)=l^2+O(l).
\]
It follows that
\(\int(\sin^{-2}\theta-1)u_l^2\,d\mu=O(l^{-1})\).
The integrand factor is bounded below by a positive constant away
from every fixed equatorial neighborhood.  The round probability
masses \(u_l^2d\mu\) therefore concentrate at the equator.  Since
\(\psi=1\) there and is continuous,
\[
 \int\psi u_l^2\,d\mu=1-o(1),\qquad
 M_l=1-\frac{2c_0}{l}+o(l^{-1}).
\]
The lowest round eigenvalue in the sector is \(l(l+1)\), so also
\(E_l\geq l(l+1)\).  The chosen high value satisfies
\begin{equation}\label{n:eq-high-parity}
 \lambda_{\rm high}=\frac{E_l}{M_l}
 \geq l(l+1)\bigl(1+2c_0/l+o(l^{-1})\bigr)
 =l^2+(1+2c_0)l+o(l).
\end{equation}
Its bounds place it in the degree-\(l\) group of the parity-\(P\)
list.

For the low value use the axisymmetric eigenfunction in the slot
occupied at round by degree \(l+1\), namely slot \(l+2\) when this
list is numbered from \(1\).  Lemma~\ref{n:lem-pole-modes} shows that its eigenvalue is simple
throughout the reflection-symmetric deformation and that its reflection
parity remains $(-1)^{l+1}=-P$.

Set
\[
 R=\frac1\pi\int_0^\pi\zeta_l(\theta)\,d\theta
   =1-\frac{c_0\bar\psi}{l},\qquad
 x(\theta)=\frac1R\int_0^\theta\zeta_l(t)\,dt .
\]
Then \(x\) maps \([0,\pi]\) onto itself, and in this coordinate
\[
 \widehat h_l=R^2\bigl(dx^2+\sin^2x\,F(x)^2d\varphi^2\bigr),
 \qquad
 F(x)=\frac{\zeta_l(\theta(x))\sin\theta(x)}{R\sin x}.
\]
The function \(F\) is positive and reflection-symmetric.  It has
smooth even extensions at \(0,\pi\) and satisfies
\(F=1+O(l^{-1})\) in every fixed smooth norm.  For example,
\(\zeta_l=1\) near a pole, so there \(\theta=Rx\) and
\(F=\sin(Rx)/(R\sin x)\), which has the stated even extension and
bounds.  The other pole follows by reflection.

The axisymmetric nonnegative Laplacian of the metric in parentheses
is
\[
 -\frac1{\sin x\,F}\partial_x(\sin x\,F\,\partial_x)
\]
on \(L^2(\sin x\,F\,dx)\).  We use the Friedrichs realizations of
the Dirichlet forms defined first on smooth profiles with even smooth
extensions at both poles.  Multiplication by \(\sqrt F\) maps this
core bijectively to the corresponding round core, since \(\sqrt F\)
and its reciprocal are smooth and even at the poles.  On this core,
the boundary contribution
\(\bigl[\sin x\,(\sqrt F)'(\sqrt F)^{-1}|v|^2\bigr]_0^\pi\)
vanishes.  The conjugated form is therefore the round axisymmetric
form on \(L^2(\sin x\,dx)\) plus the potential
\[
 V_F=\frac{(\sqrt F)''+\cot x\,(\sqrt F)'}{\sqrt F}.
\]
The even endpoint extensions make the apparent endpoint singularities
removable, and \(\|V_F\|_\infty=O(l^{-1})\).  This bounded potential
makes the shifted form norms equivalent, so multiplication by
\(\sqrt F\) also maps the closed form domain onto the round form
domain.  Min--max for a bounded potential, at the degree-\((l+1)\)
slot, gives
\begin{equation}\label{n:eq-low-parity}
 \begin{aligned}
 \lambda_{\rm low}
  &=R^{-2}\bigl((l+1)(l+2)+O(l^{-1})\bigr)\\
  &=l^2+(3+2c_0\bar\psi)l+O(1).
 \end{aligned}
\end{equation}
The same form comparison on the axisymmetric subspace places this
slot in the interval \eqref{n:eq-parity-form-bounds} for \(q=l+1\).
The separated parity intervals therefore locate it in the
degree-\((l+1)\) group of the parity-\(-P\) list.  Combining
\eqref{n:eq-high-parity} and \eqref{n:eq-low-parity} yields the
strict inversion
\begin{equation}\label{n:eq-parity-inversion}
 \lambda_{\rm high}-\lambda_{\rm low}
 \geq\bigl(2c_0(1-\bar\psi)-2\bigr)l+o(l)
 =(1-3\bar\psi)l+o(l)>0.
\end{equation}

To locate the resulting double in the combined list, write
\(\lambda_j^P\) and \(\lambda_j^{-P}\) for the ordered parity
eigenvalues.  The counts up to round degree \(l\) are
\begin{equation}\label{n:eq-parity-counts}
 A_l=\frac{(l+1)(l+2)}2\quad\text{in parity }P,\qquad
 D_l=\frac{l(l+1)}2\quad\text{in parity }-P,\qquad
 A_l+D_l=B_l.
\end{equation}
At the round metric,
\[
 \lambda_{A_l}^P=l(l+1)
 <(l+1)(l+2)=\lambda_{D_l+1}^{-P}.
\]
At \(\widehat h_l\), the high value lies among the first \(A_l\)
values of parity \(P\), while the low value lies in the group
starting at \(D_l+1\) in parity \(-P\).  Hence
\[
 \lambda_{A_l}^P\geq\lambda_{\rm high}
 >\lambda_{\rm low}\geq\lambda_{D_l+1}^{-P}.
\]
These strict indexed inequalities survive area normalization and
sufficiently small perturbations of the two endpoints.  By
Lemma~\ref{n:lem-avoidance}, choose such endpoints that are simple
through the needed cutoffs within each parity, and join them by an
even path with that property.  Along it the continuous function
\(\lambda_{A_l}^P-\lambda_{D_l+1}^{-P}\) changes sign.  At a zero,
the two equal values are simple in their separate parity lists.
Exactly
\((A_l-1)+D_l=B_l-1\) eigenvalues of the combined list are below
them.  Thus the equality is a double at positions \(B_l,B_l+1\).
This conclusion requires only continuity of the indexed parity
eigenvalues; no isolation of coincidences along this preliminary
path is needed.  Apply Lemma~\ref{n:lem-retained-double} in
unrestricted directions to keep this double and separate every
other level through \(p+1\).

The two cases cover all adjacent pairs in \(I\).  For fixed \(k\)
there are finitely many of them.  All cross-block estimates above are
estimates as \(l\to\infty\), so they hold simultaneously for
\(L+1\leq l\leq M-1\) once \(L\) is large.  The cross-block parameters
\(w_l=\log(1-c_0\psi/l)\) tend smoothly to zero uniformly for
\(l\geq L+1\) as \(k\) tends to infinity; the remaining
perturbations can be arbitrarily small.  Thus one prescribed
\(\mathcal U\) contains all the metrics selected for every
sufficiently large \(k\).
\end{proof}
